\markboth{Erstes Kapitel}{Erstes Kapitel}
\chapter{Kapitel mit Testelementen - LaTeX}                                         % 1     Kapitel mit Testelementen - LaTeX
\section{Zitierungen}                                                               % 1.1   Zitierungen
Das Zitieren im Harvard-Stil ist vor allem bei Sozial- und
Naturwissenschaftlern beliebt. Ein Hauptunterschied zu anderen Zitierweisen
ist dabei der Verzicht auf Fu\ss noten. Hei\ss t, dass du mitunter mehr Platz
pro Seite für deinen Flie\ss text hast - ob das nun ein Fluch oder Segen ist,
\"uberlassen wir deiner Vorliebe. Der Verzicht auf Fu\ss noten im
Harvard-Zitieren zieht nat\"urlich eines nach sich: Den Literaturhinweis
musst du in deinen Text einflie\ss en lassen. Doch auch da zieht die
amerikanische Elite-Universit\"at die Sparsamkeit der Vollst\"andigkeit vor.
An geeigneter Stelle verweist du n\"amlich nur auf den Nachnamen der Autorin
oder des Autors, das Erscheinungsjahr und die Seite, von der du zitierst.\par

EIn Zitat kann dabei durch eine hochgestellte Fu\ss note
\footfullcite{aristotle:physics} oder durch einen Literaturhinweis im Sinne der
Harvard-Zitierweise.                                                                % footfullcite (ausführlich) oder \footcite (kurz)

\subsection{Direktes Zitat}                                                         % 1.1.1 Direktes Zitat
Die Funktion der Ehe ist als
"Grundstein der Familie und als Zentrum der geistig-seelischen Entwicklung des
Menschen und andererseits als \"okonomisch-politischer Nutzfaktor"
\autocite{aristotle:physics} zu betrachten.                                         % Variante (ohne Fu\ss zeile, Name, Jahr)

\subsection{Indirektes Zitat}
Die Architextualit\"at besch\"aftigt sich damit, wie die Zuordnung eines Textes
zu einer Kategorie stattfindet und einer Gattung zugeordnet wird
\autocite[9]{aristotle:physics}.                                                    % Variante (ohne Fu\ss zeile, Name, Jahr, Seite)

\subsection{Auslassungen}
Die Funktion der Ehe ist als "Grundstein der Familie und als Zentrum der
geistig-seelischen Entwicklung des Menschen [...]"
\autocite[9]{aristotle:physics} zu betrachten.                                      % Variante (ohne Fu\ss zeile, Name, Jahr, Seite)

\subsection{Ver\"anderungen}
Die Funktionen von Ehen sind die "Grundstein[e] der Familie und als Zentrum
der geistig-seelischen Entwicklung des Menschen und andererseits als
\"okonomisch-politischer Nutzfaktor" \autocite[9]{aristotle:physics} zu betrachten. % Variante (ohne Fu\ss zeile, Name, Jahr, Seite)

\subsection{Falschschreibungen}
Befinden sich Tipp-, Grammatik- oder \"ahnliche Fehler im Original,
\"ubernimmst du diese zwar, kennzeichnest dies jedoch durch ein [sic!]
hinter dem jeweils fehlerhaften Wort. Aus Gr\"unden der Lesbarkeit
wird auf ein [sic!] verzichtet, wenn Textstellen zitiert werden,
die noch in alter Rechtschreibung geschrieben sind.                                 % Fehler im Orginal

\section{Indexeintr\"age}
Affe\index{Affe}
Bieber\index{Bieber}                                                                % Indexverweis durch \index{Bez}


\section{Abk\"urzungen}                                                             
Das K\"urzel EU \ac{EU} wird in erster Verwendung auch in der Langform
dargestellt. Wird das gleiche K\"urzel (\ac{EU}) erneut aufgerufen,
so erfolgt lediglich die Kurzform.                                                  % Abk\"urzungen durch \ac{Bez}

\section{Abbildungen}
\begin{figure}
    \centering                                                                      % zentrieren der Abbildung
    \includegraphics[width=.9\linewidth]{ba_abbildungen/testimage}                  % .ps, width=10cm, scale=0.5, width=.5\linewidth
    \captionof{figure}{Eine Beispielabbildung}                                      % \Beschriftung der Abbildung
\end{figure}

\section{Tabellen}
\begin{table}
    \centering                                                                      % zentrieren der Tabelle
    \begin{tabular}{l|c|r}\hline\hline                                              % Definition der Tabelle (spalten, links-, zentriert- oder rechts-bündig)
        links & zentriert & rechts \\\hline A & B & B\\\cline{2-2}                  % Verbinden von Zeilen
        1 & 2 & 3\\\hline
    \end{tabular}
    \captionof{table}{Tabellenbeschreibung}                                         % Beschriftung der Tabelle
\end{table} 

\newpage